\documentclass[11pt]{article}
\usepackage[utf8]{inputenc}
\usepackage{amsmath,amssymb}
\usepackage{graphicx}
\usepackage{booktabs}
\usepackage{authblk}
\usepackage[margin=1in]{geometry}
\usepackage[hidelinks]{hyperref}
\usepackage{xurl}
\usepackage{setspace}

\title{Complex Time-Dependent RealQM:\\
Currents, Radiation, and Magnetism from the Real-Space Continuum}
\author[1]{Claes Johnson}
\affil[1]{KTH Royal Institute of Technology, SE-100 44 Stockholm, Sweden \\
\texttt{claesjohnson@gmail.com}}
\date{\today \\[0.4em]
\small\textit{Argument and formulation by the author; drafting and computation prepared with AI assistance (Claude, Anthropic), reviewed and verified by the author.}}

\begin{document}
\maketitle

\begin{abstract}
RealQM in its established form is a \emph{real-valued}, \emph{parabolic} theory: electrons are
real-valued charge densities $\psi_i:\mathbb R^3\to\mathbb R$ on non-overlapping domains, and ground
states are found by imaginary-time (gradient-flow) relaxation of the Coulomb energy. This form is
validated across chemistry, but it produces no electric \emph{current} --- a real wavefunction has
$\mathbf J=\frac{\hbar}{m}\mathrm{Im}(\psi^*\nabla\psi)\equiv 0$ --- and therefore no radiation currents,
no magnetic moments, and no spin observables; the base theory itself flags these as requiring a separate
extension. Here we formulate that extension: a \emph{complex}, \emph{unitary}, time-dependent RealQM,
$i\hbar\,\partial_t\psi_i = \mathcal H_i\psi_i$ on each domain, carrying a well-defined charge current,
alongside the existing parabolic relaxation. We show that the free-boundary (Bernoulli / zero-flux)
condition that already defines the domains \emph{automatically} makes the unitary evolution
charge-conserving on each domain, so the extension respects the multiphase structure with no new
assumption. Two kinds of motion then generate the electromagnetic observables the real theory cannot:
\emph{oscillation} (a superposition of states, a time-varying density) sources dipole radiation, and
\emph{circulation} (a spatial phase winding, a standing current) gives a magnetic moment
$\boldsymbol\mu=\tfrac12\int\mathbf r\times\mathbf J$. Ground states are real, hence currentless, hence
carry no moment --- which is why RealQM's flat confined electron has none --- while excited or winding
states carry both. We are explicit that this is a proposed formulation, consistent with but extending the
validated real theory: the ground-state energetics are unchanged; the current-bearing dynamics, and its
use for the free electron's own moment and for spin, are new and open.

\medskip
\noindent\textbf{Keywords:} RealQM; time-dependent; complex wavefunction; charge current; free boundary;
radiation; magnetic moment; spin; parabolic vs unitary.
\end{abstract}

\setstretch{1.12}

\section{Why an extension is needed}
\label{sec:why}

RealQM \cite{RealQM} represents the $N$ electrons of an atom or molecule by real-valued one-electron wave
functions $\psi_i:\mathbb R^3\to\mathbb R$ supported on non-overlapping domains $\Omega_i$ with free
(Bernoulli) boundaries, and finds the ground state by minimising the Coulomb energy
\begin{equation}
E \;=\; \sum_i \tfrac12\!\int_{\Omega_i}|\nabla\psi_i|^2
\;+\; \tfrac12\!\iint\frac{\rho(\mathbf x)\rho(\mathbf y)}{|\mathbf x-\mathbf y|}\,d\mathbf x\,d\mathbf y,
\qquad \rho=\sum_i q_i\,\psi_i^2 ,
\label{eq:E}
\end{equation}
via imaginary-time relaxation $\dot\psi_i + \mathcal H_i\psi_i = 0$, a structurally \emph{parabolic}
(dissipative) flow toward the lowest state. This is the validated core, and it is deliberately real.

A real wavefunction, however, carries \emph{no current}: with $\psi_i$ real,
$\mathrm{Im}(\psi_i^*\nabla\psi_i)=0$. So the base theory cannot express any observable that is built from
charge \emph{in motion} --- dipole radiation, magnetic moments, magnetic susceptibility, NMR shifts,
hyperfine structure. RealQM says as much of itself, flagging these as outside its current form and
requiring ``a separate extension.'' This note writes that extension down.

The guiding principle is minimal: keep the domains, the Coulomb energy \eqref{eq:E}, and the free-boundary
condition exactly as they are, and add only what is needed for motion --- a \emph{complex} phase and a
\emph{unitary} time law that carries it.

\section{The complex time-dependent formulation}
\label{sec:form}

In the formulation below we set $\hbar=1$; the masses $m_i$ and charges $q_i=\pm1$ are kept explicit, since different
domains may carry different mass (electrons, protons, or the equal-mass device of the nuclear regime),
and the current depends on that mass.

Let each electron wave function now be complex and time-dependent,
$\psi_i:\Omega_i\times\mathbb R\to\mathbb C$, with charge density and charge current
\begin{equation}
\rho_i = q_i\,|\psi_i|^2, \qquad
\mathbf J_i = \frac{q_i}{m_i}\,\mathrm{Im}\!\left(\psi_i^*\nabla\psi_i\right),
\qquad \rho=\sum_i\rho_i,\ \ \mathbf J=\sum_i\mathbf J_i .
\label{eq:rhoJ}
\end{equation}
The current \eqref{eq:rhoJ} has a transparent hydrodynamic (Madelung) reading. Writing
$\psi_i = R_i\,e^{iS_i}$ with real amplitude $R_i$ and real phase $S_i$, one has
$\psi_i^*\nabla\psi_i = R_i\nabla R_i + i\,R_i^2\,\nabla S_i$, so
$\mathrm{Im}(\psi_i^*\nabla\psi_i) = |\psi_i|^2\,\nabla S_i$ and
\begin{equation}
\mathbf J_i \;=\; \rho_i\,\mathbf v_i, \qquad \mathbf v_i = \frac{\nabla S_i}{m_i} .
\label{eq:madelung}
\end{equation}
The current is charge density times a \emph{velocity field} $\mathbf v_i$ set by the phase gradient
($\mathbf p_i = \nabla S_i$ is the local momentum, $\mathbf v_i = \mathbf p_i/m_i$). Two facts follow at
once: a \emph{real} density has constant phase, $\nabla S_i = 0$, hence $\mathbf v_i = 0$ and
$\mathbf J_i = 0$ --- no motion, no current; and a phase that \emph{winds} in space, $S_i = m\phi$, gives
an azimuthal velocity $|\mathbf v_i| = m/(m_i s)$ and a standing circulation. The mass enters only through
$\mathbf v_i = \nabla S_i/m_i$: it scales the current when there \emph{is} a phase gradient, and is
irrelevant when there is none --- consistent with the flat electron being both mass-irrelevant and
currentless.

The evolution is \emph{unitary} on each domain,
\begin{equation}
i\,\partial_t\psi_i \;=\; \mathcal H_i\,\psi_i,
\qquad
\mathcal H_i = -\frac{1}{2m_i}\nabla^2 + V_i ,
\label{eq:tdse}
\end{equation}
with $V_i$ the same real (kernel plus inter-domain Hartree) potential as in the ground-state problem,
$V_i(\mathbf x)=W(\mathbf x)+\sum_{j\ne i}\int q_j|\psi_j|^2/|\mathbf x-\mathbf y|\,d\mathbf y$. The free
boundary carries the \emph{same} Bernoulli condition as before,
\begin{equation}
\frac{\partial\psi_i}{\partial n}=0 \quad\text{on }\Gamma_i=\partial\Omega_i .
\label{eq:neumann}
\end{equation}

\paragraph{The free boundary conserves charge automatically.} Equations \eqref{eq:tdse}--\eqref{eq:neumann}
have a property that makes the extension consistent with the multiphase structure without any new
assumption. From \eqref{eq:tdse} with real $V_i$ follows the continuity equation
\begin{equation}
\partial_t\rho_i + \nabla\!\cdot\!\mathbf J_i = 0 \quad\text{in }\Omega_i .
\label{eq:cont}
\end{equation}
Integrating over $\Omega_i$ and using the divergence theorem, the rate of change of the charge on domain
$i$ is the flux through its boundary,
$\dot Q_i = -\oint_{\Gamma_i}\mathbf J_i\!\cdot\!\mathbf n\,dS$. But the normal current is
$\mathbf J_i\!\cdot\!\mathbf n = (q_i/m_i)\,\mathrm{Im}(\psi_i^*\,\partial_n\psi_i)$, and the Bernoulli
condition \eqref{eq:neumann} sets $\partial_n\psi_i=0$, so
\begin{equation}
\mathbf J_i\!\cdot\!\mathbf n = 0 \ \text{on }\Gamma_i
\qquad\Longrightarrow\qquad \dot Q_i = 0 .
\label{eq:noflux}
\end{equation}
\textbf{Each electron keeps its unit charge under the unitary evolution}, because the very boundary
condition that defines the RealQM domains is also a zero-current (no-leak) condition. The free boundary,
introduced for the statics, is exactly what the complex dynamics needs to remain a clean $N$-electron
partition.

\paragraph{Are the domains held fixed?} In the minimal formulation adopted here, yes: the domains
$\Omega_i$ are kept at their ground-state shape and each $\psi_i$ evolves \emph{within} its fixed domain
under \eqref{eq:tdse}, subject to the zero-flux boundary condition \eqref{eq:neumann}--\eqref{eq:noflux}
that conserves its charge. This is exact for small-amplitude motion, and --- importantly for the
applications here --- it is well-justified for a \emph{winding} state, because a winding
$\psi_i = R_i\,e^{im\phi}$ changes only the \emph{phase} of the density, not its magnitude: $|\psi_i|^2 =
R_i^2$ is essentially the ground-state density (with, at most, a node forced onto the winding axis), so
the free boundary, which is fixed by $|\psi_i|^2$, barely moves. The magnetic moment is exactly such a
winding state, so holding the domains fixed is not a crude approximation but a controlled one. The
general, large-amplitude case --- in which the free boundary $\Gamma_i$ co-evolves with the $\psi_i$, the
domains genuinely moving --- is the harder problem, and its well-posedness is listed as open
(\S\ref{sec:open}). What must not be lost is that in RealQM the partition is normally a \emph{primary}
variational degree of freedom; fixing it is a deliberate small-amplitude simplification, exact at rest and
controlled for windings, not a claim that the domains are permanently rigid.

\paragraph{Relation to the parabolic form.} The two dynamics are the same operator under
$\tau=it$: the parabolic relaxation $\dot\psi_i=-\mathcal H_i\psi_i$ drives $\psi_i$ to the real ground
state (imaginary time), while the unitary law \eqref{eq:tdse} carries the complex phase (real time).
\emph{Relaxation finds the state; unitary evolution carries the current.} A system prepared in its ground
state and left alone has real $\psi_i$, hence $\mathbf J=0$: no radiation, no moment, at rest --- as it
should be.

\section{The electromagnetic observables the real theory could not give}
\label{sec:obs}

With a current in hand, the observables follow by the classical (Schr\"odinger--Maxwell) reading: the
charge density and current $\rho,\mathbf J$ source the electromagnetic field, and every electromagnetic
observable is a functional of them. Two distinct motions matter.

\paragraph{Oscillation $\Rightarrow$ radiation.} A superposition of two stationary states,
$\psi_i = a\,\phi_1 e^{-iE_1t}+b\,\phi_2 e^{-iE_2t}$, has a density that beats at
$\omega=E_2-E_1$ and an oscillating dipole $\mathbf d(t)=\int\mathbf x\,\rho\,d\mathbf x$, which
radiates at $\omega$ by the Larmor formula. This is the charge ``oscillating between two states'' of the
matter--radiation picture, now with a definite current behind it.

\paragraph{Circulation $\Rightarrow$ magnetic moment.} A stationary state with a \emph{spatial phase
winding}, $\psi_i = R(\mathbf x)\,e^{im\phi}$ about an axis, carries a standing azimuthal current and a
magnetic moment
\begin{equation}
\boldsymbol\mu \;=\; \tfrac12\!\int \mathbf r\times\mathbf J\,d^3x ,
\qquad \mu_z = -\,m\,\mu_B \quad(\text{for a normalised electron density}),
\label{eq:mu}
\end{equation}
the moment being quantised in the winding number $m$ and independent of the density's shape. A real
($m=0$) density gives $\mu=0$. This is the magnetostatic partner of the oscillation: charge going in
circles rather than back and forth.

\paragraph{Why the confined electron has no moment.} The two facts combine into the result of the
companion note \cite{NuclearMoment}: RealQM computes the deuteron's caged electron as a \emph{flat, real}
density --- the $m=0$, currentless case --- so it carries no magnetic moment, and nuclear moments come out
at nuclear-magneton scale. That result lives already in the \emph{real} theory ($\psi$ real $\Rightarrow$
$\mathbf J=0$); the present extension is what would be needed for the \emph{nonzero} moments (the free
electron's own $\mu_B$, the neutron's residual), which require a genuine current.

\section{Spin as a further structure}
\label{sec:spin}

Spin is not a primary variable of RealQM; Pauli exclusion is realised geometrically by non-overlap. The
natural place for spin in the complex extension is an \emph{internal circulation}: the free electron's
$\mu_B$ cannot be an ordinary orbital winding (the free/atomic ground state is real, $m=0$, with zero
orbital moment), so if it is a current at all it is a Compton-scale internal circulation --- the
Zitterbewegung structure --- carrying angular momentum $\tfrac12$ and moment $\sim\mu_B$. Equivalently, one
may add the $\pm\tfrac12$ domain label the base theory suggests, with its flux balance entering the
magnetic field. Either route is a genuine addition beyond \eqref{eq:tdse}; we flag it, and the anomalous
$g/2=1.00116\ldots$, as open (\S\ref{sec:open}) rather than claim them here.

\section{What is proposed, what is validated, what is open}
\label{sec:open}

\emph{Validated (unchanged):} the real, parabolic ground-state RealQM and all its chemistry --- energies,
geometries, reactions, condensed phases --- are untouched. The extension reduces to it exactly for a
system at rest (real $\psi$, $\mathbf J=0$).

\emph{Proposed here:} the complex, unitary, time-dependent law \eqref{eq:tdse} on the RealQM domains, its
current \eqref{eq:rhoJ}, and the observation \eqref{eq:noflux} that the Bernoulli free boundary makes it
charge-conserving per domain. This gives radiation and (orbital) magnetic moments a first-principles home
in the real-space continuum.

\emph{Open:}
\begin{itemize}
\item \textbf{The moving free boundary.} For large-amplitude dynamics the domains $\Omega_i$ should
co-evolve with the $\psi_i$; the well-posedness of the coupled complex-field / free-boundary problem is a
genuine mathematical question. The fixed-domain version is the clean small-amplitude case.
\item \textbf{The free electron's $\mu_B$ and spin.} These require the Compton-scale internal circulation
of \S\ref{sec:spin}, not the atomic-scale winding \eqref{eq:mu}; they are not delivered by
\eqref{eq:tdse} alone.
\item \textbf{The anomalous moment.} $g/2 = 1.00116\ldots$ (QED's most precise number) would be the sharp
test of any current-based account of the electron moment, and is untouched here.
\item \textbf{Numerics.} Implementing \eqref{eq:tdse} (two real fields per domain, a norm-preserving
staggered-leapfrog step, the current \eqref{eq:rhoJ}) on top of the existing solver is the computational
next step, and the vehicle for the free-vs-caged moment test.
\end{itemize}

\section{Conclusion}
\label{sec:concl}

RealQM was built real and parabolic, and is validated there; but charge in motion --- radiation,
magnetism, spin --- lives only in a complex, unitary, time-dependent form, which the base theory itself
marks as missing. The extension is minimal: the same domains, the same Coulomb energy, the same Bernoulli
boundary, now with a complex phase evolved by $i\,\dot\psi_i=\mathcal H_i\psi_i$. The free boundary
turns out to be exactly the zero-current condition that keeps the evolution a clean $N$-electron
partition, so nothing new must be assumed. Oscillation then gives radiation and circulation gives a
magnetic moment, while the real ground state --- currentless --- gives neither, which is why the flat
confined electron carries no moment. What remains open, honestly, is the free electron's own $\mu_B$ at
the Compton scale, the anomalous $g$-factor, the moving-boundary dynamics, and the numerics --- but the
formulation on which all of these can now be posed is in place.

\begin{thebibliography}{9}
\bibitem{RealQM} C.~Johnson, \emph{Real Quantum Mechanics as Multiphase 3D Continuum Mechanics},
manuscript, 2026; \url{https://claes542.github.io/RealMolecule/gallery.html}.
\bibitem{NuclearMoment} C.~Johnson, \emph{Magnetism as Circulating Charge: Why the Nuclear Electron
Carries No Magnetic Moment}, 2026 (companion note).
\end{thebibliography}

\end{document}
